\documentclass[11pt,a4paper]{article}
\usepackage{qstyle}
\begin{document}

% ----------------------------------------------------------------- cover
\thispagestyle{empty}
\vspace*{6mm}
{\color{inkred}\Huge\bfseries BAV --- Exam Question Bank}\par\vspace{3mm}
{\color{inkblue}\Large Business Analysis \& Valuation}\par\vspace{1mm}
{\color{pencil}\large 45 solved multiple-choice questions}\par\vspace{6mm}

\begin{minipage}{0.97\textwidth}
{\color{inkblue}
Forty questions on the taught course, plus five on the guest-lecture session.
Every question carries four options, a full worked solution and the answer. Numbers
are taken from the Ambuja and UltraTech valuations built for the project, so the
arithmetic is checkable against the models.}
\end{minipage}

\vspace{7mm}
\begin{minipage}{0.97\textwidth}
{\color{inkblue}\bfseries What is in here}\par\vspace{2mm}
\begin{tabular}{@{}p{0.13\textwidth}p{0.56\textwidth}p{0.11\textwidth}@{}}
\textbf{A} & The five consistencies & Q1--8\\[1.5mm]
\textbf{B} & DCF variants and why they must agree & Q9--17\\[1.5mm]
\textbf{C} & Cost of capital, relevering and circularity & Q18--26\\[1.5mm]
\textbf{D} & Terminal value, multiples and relative valuation & Q27--35\\[1.5mm]
\textbf{E} & Applied: the cement valuations & Q36--40\\[1.5mm]
\textbf{F} & Guest lecture: M\&A, synergy and exchange ratios & Q41--45\\
\end{tabular}
\end{minipage}

\vspace{8mm}
\begin{minipage}{0.97\textwidth}
{\color{inkred}\bfseries How to use it.} {\color{inkblue}Cover the working with a
sheet of paper, answer the four options, then check. The starred
\textbf{$\star$ Remember} lines are the traps that separate marks --- most of them are
sign errors, wrong-rate errors or double counting, not hard arithmetic.}
\end{minipage}

\vspace{6mm}
\begin{minipage}{0.97\textwidth}
{\color{pencil}\itshape Note on Section F: the guest-lecture questions are written on
M\&A, synergy valuation and exchange ratios, which is the module the course carries
separately from the core valuation syllabus. If the guest session covered a different
topic, these five should be swapped.}
\end{minipage}

\newpage

% ================================================================= SECTION A
\sect{A}{The five consistencies}{Mohanty's central claim: a valuation is correct when
it preserves five consistencies --- assumptions, risk, asset, liability, multiplier.
Maintain all five and every DCF variant returns the same equity value.}

\q{1}{Five consistencies --- risk}{A student discounts free cash flow to \emph{equity}
at the WACC. Which consistency has been broken?}
\opt{Assumptions consistency}{Risk consistency}{Asset consistency}{Multiplier consistency}
\begin{soln}
Risk consistency requires the cash flow and the discount rate to belong to the same
claimholder. FCFE is an \emph{equity} flow, so it must be discounted at the cost of
equity $K_e$. WACC is a firm-level rate and belongs with FCFF or CCF.
\formula{Firm flows $\to$ WACC or $\rho$ \quad\textbar\quad Equity flows $\to$ $K_e$}
Since WACC $<K_e$ whenever there is debt, this error \emph{overstates} equity value.
\end{soln}
\ans{b}{Risk consistency}
\memo{The classic pair is FCFE at WACC and FCFF at $K_e$. Check the label on the cash
flow before you check the arithmetic.}

\q{2}{Five consistencies --- assumptions}{A forecast grows sales at 15\% a year while
holding capital expenditure and working capital flat in rupee terms. Which consistency
fails?}
\opt{Assumptions consistency}{Liability consistency}{Risk consistency}{Multiplier consistency}
\begin{soln}
Growth requires reinvestment. Selling more cement needs more plant and more working
capital; a model that grows revenue without growing the capital base is manufacturing
free cash flow out of nothing.
\formula{$g = \text{RONIC} \times \text{reinvestment rate}$}
Rearranged, reinvestment rate $=g/\text{RONIC}$. If reinvestment is zero, $g$ must be
zero unless RONIC is infinite.
\end{soln}
\ans{a}{Assumptions consistency}
\memo{Every growth assumption has a matching investment assumption. If you cannot show
where the money went, the growth is not real.}

\q{3}{Five consistencies --- asset}{A model excludes interest income from NOPAT because
cash is treated as non-operating, but forgets to add the cash balance back in the
equity bridge. What is the effect?}
\opt{Equity value is unchanged}{Equity value is overstated}{Equity value is
understated}{Only the WACC is affected}
\begin{soln}
Asset consistency: every asset is counted exactly \emph{once}. Two valid treatments
exist ---
\formula{(i) income \emph{in} NOPAT $\to$ asset \emph{not} added back}
\formula{(ii) income \emph{excluded} $\to$ asset added back at the bridge}
Excluding the income and also omitting the add-back counts the cash \emph{zero} times.
The company owns the cash but the valuation gives no credit for it, so value is
understated by the cash balance.
\end{soln}
\ans{c}{Equity value is understated}
\memo{Asset errors come in both directions. Adding back land while its rent still sits
in operating income counts it twice; this question is the mirror image.}

\q{4}{Five consistencies --- liability}{A deferred sales-tax liability is included as a
source of funds in the WACC weights but is not subtracted in the bridge from enterprise
value to equity. Which consistency fails, and in which direction?}
\opt{Liability; equity overstated}{Liability; equity understated}{Asset; equity
overstated}{Risk; no directional effect}
\begin{soln}
Liability consistency: whatever is included in the cost of capital as a source of funds
\emph{must} be subtracted from enterprise value to reach equity. Including the
liability in the weights but not in the bridge means the claim finances the assets but
never gets repaid out of the value.
\formula{Equity $=$ EV $+$ non-operating assets $-$ \emph{every} prior claim}
The omitted deduction inflates the residual left for shareholders.
\end{soln}
\ans{a}{Liability; equity overstated}
\memo{Non-interest-bearing items are either inside net working capital or treated as
debt-like --- never both, and never neither.}

\q{5}{Five consistencies --- multiplier}{A model assumes terminal growth of 10\% and
produces an implied terminal EV/EBITDA of 4x. What does this indicate?}
\opt{Nothing; multiples and DCF are unrelated}{Multiplier inconsistency}{The WACC must
be too low}{The forecast horizon is too long}
\begin{soln}
A multiple is the \emph{output} of value divided by an accounting indicator, so the
multiple your model implies must be arithmetically consistent with the growth, return
and discount rate the model itself uses.
\formula{Justified $\dfrac{EV}{EBITDA} = \dfrac{(NOPAT/EBITDA)\,(1-g/ROIC)}{WACC-g}$}
High terminal growth with a very low exit multiple is internally contradictory: a
company growing at 10\% forever would command a high multiple, not 4x.
\end{soln}
\ans{b}{Multiplier inconsistency}

\q{6}{Method equivalence}{All five DCF methods in a model return exactly the same value
per share. What does this prove?}
\opt{The assumptions are correct}{The valuation is accurate}{The model is internally
consistent, and nothing more}{The market price must be wrong}
\begin{soln}
The five methods are one model rearranged. Agreement demonstrates that no consistency
has been broken in the \emph{construction} of the model. It says nothing whatever about
whether the inputs are right.
\formula{Agreement $\Rightarrow$ internal consistency \quad\textbf{but}\quad
internal consistency $\nRightarrow$ correct value}
A model with a badly wrong growth assumption will still reconcile across all five
methods, perfectly and pointlessly.
\end{soln}
\ans{c}{The model is internally consistent, and nothing more}
\memo{This is a favourite one-mark question. Never write that reconciliation proves the
valuation is right.}

\q{7}{Five consistencies --- liability}{Operating leases are capitalised and the lease
liability is deducted in the equity bridge, but rent continues to be expensed inside
the free cash flow. What has happened?}
\opt{Correct treatment}{The lease is counted twice as a cost}{The lease is not counted
at all}{Only the beta is affected}
\begin{soln}
Capitalising the lease means treating it as debt: the liability is deducted at the
bridge and the rent should be \emph{removed} from operating expense (replaced by
depreciation and interest, both of which sit below FCFF).
Leaving rent in the cash flow \emph{and} deducting the liability charges the company
for the same obligation twice --- once through lower FCFF and once through a smaller
bridge.
\end{soln}
\ans{b}{The lease is counted twice as a cost}

\q{8}{Terminal reinvestment}{A model charges terminal reinvestment as $g$ times invested
capital \emph{including} acquisition goodwill. What does this implicitly assume?}
\opt{That the company must keep buying goodwill forever merely to grow}{That goodwill
depreciates}{That the tax shield is perpetual}{Nothing --- it is the standard treatment}
\begin{soln}
Reinvestment in the terminal year funds the assets that produce the extra output.
Organic growth needs more plant and more working capital; it does not need the company
to keep making acquisitions.
\formula{Terminal reinvestment $= g \times$ \emph{operating} invested capital}
Goodwill still belongs \emph{inside} invested capital when measuring the return
shareholders earned on what they paid --- but not inside the reinvestment charge. In
the Ambuja model, getting this wrong drove the equity value negative.
\end{soln}
\ans{a}{That the company must keep buying goodwill forever merely to grow}
\memo{Goodwill: in the ROIC denominator, out of terminal reinvestment.}

% ================================================================= SECTION B
\sect{B}{DCF variants and why they must agree}{The OLP Limited test vector runs through
Q9--Q12. Inputs: net operating income $X=\text{Rs }100$m in perpetuity, tax
$t=20\%$, debt $D=\text{Rs }100$m at $K_d=10\%$, cost of equity $K_e=20\%$, no
non-operating assets, zero growth.}

\q{9}{OLP --- FCF at WACC}{Using the OLP inputs, what is the enterprise value under the
enterprise FCF method?}
\opt{Rs 400m}{Rs 440m}{Rs 460m}{Rs 500m}
\begin{soln}
Step 1 --- the cash flow. With zero growth there is no reinvestment, so
$\text{FCF} = X(1-t) = 100(0.8) = \text{Rs }80$m.\par
Step 2 --- equity value, needed for the weights. $\text{FCFE}=(X-K_dD)(1-t)
=(100-10)(0.8)=72$, so $E = 72/0.20 = \text{Rs }360$m and $V = 360+100 = 460$.\par
Step 3 --- the WACC.
\formula{$\text{WACC}=\tfrac{E}{V}K_e+\tfrac{D}{V}K_d(1-t)
=\tfrac{360}{460}(0.20)+\tfrac{100}{460}(0.10)(0.8)=17.391\%$}
Step 4 --- $EV = 80/0.17391 = \textbf{Rs }\mathbf{460}$m.
\end{soln}
\ans{c}{Rs 460m}

\q{10}{OLP --- FCFE at $K_e$}{Using the same inputs, what is the equity value under the
FCFE method?}
\opt{Rs 360m}{Rs 400m}{Rs 460m}{Rs 288m}
\begin{soln}
Equity cash flow is what is left after paying interest and tax, with zero growth so no
net new borrowing.
\formula{$\text{FCFE}=(X-K_dD)(1-t)=(100-10)(0.8)=\text{Rs }72$m}
Discount at the \emph{equity} rate: $E = 72/0.20 = \textbf{Rs }\mathbf{360}$m.\par
Cross-check: $EV = E+D = 360+100 = 460$, matching Q9 exactly. That equality is the
whole point of the exercise.
\end{soln}
\ans{a}{Rs 360m}

\q{11}{OLP --- CCF at pre-tax WACC}{What is the capital cash flow, and the rate at which
it is discounted?}
\opt{80 at 17.391\%}{82 at 17.826\%}{72 at 20.000\%}{82 at 17.391\%}
\begin{soln}
Capital cash flow adds the interest tax shield back into the flow, so the shield is no
longer inside the discount rate --- which means the rate must be the \emph{pre-tax}
WACC.
\formula{$\text{CCF}=\text{FCF}+t\,K_d D = 80 + (0.20)(0.10)(100) = \text{Rs }82$m}
\formula{$\text{pre-tax WACC}=\tfrac{E}{V}K_e+\tfrac{D}{V}K_d
=\tfrac{360}{460}(0.20)+\tfrac{100}{460}(0.10)=17.826\%$}
$EV = 82/0.17826 = \text{Rs }460$m. Same answer again.
\end{soln}
\ans{b}{82 at 17.826\%}
\memo{If you add the shield to the numerator you must take it out of the denominator.
Doing both, or neither, is the error.}

\q{12}{OLP --- APV}{Under APV with debt fixed in rupee terms, what are the unlevered
cost of capital and the value of the tax shield?}
\opt{$K_u=20.000\%$, VTS $=$ Rs 10m}{$K_u=18.182\%$, VTS $=$ Rs 20m}{$K_u=17.391\%$,
VTS $=$ Rs 20m}{$K_u=18.182\%$, VTS $=$ Rs 80m}
\begin{soln}
Unlever using the MM (1963) relation, which is the fixed-debt case:
\formula{$K_e = K_u + (K_u-K_d)\tfrac{D}{E}(1-t)$}
$0.20 = K_u + (K_u-0.10)\tfrac{100}{360}(0.8)$. Solving, $K_u = \mathbf{18.182\%}$.\par
With debt fixed and zero growth the shield is a perpetuity discounted at $K_d$, which
collapses to
\formula{$\text{VTS} = \dfrac{t K_d D}{K_d} = tD = 0.20 \times 100 = \text{Rs }20$m}
$\text{APV} = 80/0.18182 + 20 = 440 + 20 = \text{Rs }460$m.
\end{soln}
\ans{b}{$K_u=18.182\%$, VTS $=$ Rs 20m}

\q{13}{FCFF to FCFE bridge}{Which identity converts free cash flow to the firm into free
cash flow to equity?}
\opt{FCFE $=$ FCFF $-$ interest $+$ tax}{FCFE $=$ FCFF $-$ interest$(1-t)$ $+$ net new
borrowing}{FCFE $=$ FCFF $+$ interest$(1-t)$ $-$ net new borrowing}{FCFE $=$ FCFF $-$
dividends}
\begin{soln}
FCFF is the cash available to \emph{all} providers of capital. To get to the equity
holders you must pay the lenders and then add back anything the lenders newly advanced.
\formula{$\text{FCFE} = \text{FCFF} - I(1-t) + \Delta D$}
Interest is taken after tax because it is tax-deductible. Net new borrowing $\Delta D$
is a genuine cash inflow to equity.
\end{soln}
\ans{b}{FCFE $=$ FCFF $-$ interest$(1-t)$ $+$ net new borrowing}

\q{14}{Net new borrowing}{A firm grows at $g$ forever and maintains a constant debt
ratio. What is net new borrowing each year?}
\opt{Zero}{$g \times D$}{$K_d \times D$}{$t \times K_d \times D$}
\begin{soln}
A constant debt \emph{ratio} with a firm growing at $g$ means the rupee amount of debt
must also grow at $g$; otherwise the ratio falls every year.
\formula{$\Delta D = g \times D$}
This is the most common error in the course: students discount a growing FCFE at a
constant $K_e$ while implicitly holding debt fixed. Equity flows then grow faster than
$g$, the debt ratio drifts, and the constant-$K_e$ assumption collapses.
\end{soln}
\ans{b}{$g \times D$}
\memo{If you hold WACC constant you have assumed a constant D/E, which \emph{requires}
debt to grow at $g$. Put the net borrowing line in the model and show it.}

\q{15}{Choice of method}{A company is deleveraging on a announced schedule over five
years. Which valuation method is most appropriate?}
\opt{FCFF at a single constant WACC}{APV}{Relative valuation only}{Dividend discount
model}
\begin{soln}
A single constant WACC embeds a constant capital structure. If leverage is changing on
a known path, the WACC changes every year and one average rate is simply wrong.\par
APV separates the two effects: value the assets unlevered at $K_u$, then add the
present value of the year-by-year tax shields the actual debt schedule produces.
\formula{$\text{APV} = V_u + \text{PV(tax shields)}$}
\end{soln}
\ans{b}{APV}
\memo{APV is the LBO and deleveraging method precisely because the financing side is
modelled explicitly rather than buried in a rate.}

\q{16}{EVA reconciliation}{Enterprise value under EVA equals opening invested capital
plus the present value of future EVA. Which two conditions are necessary for this to
tie to the FCFF answer?}
\opt{Positive EVA, and a five-year horizon}{FCFF defined as NOPAT less the change in
invested capital, and a true steady-state terminal year}{A constant WACC, and zero
growth}{Goodwill excluded from invested capital}
\begin{soln}
EVA is not an independent method; it is the same cash flows re-expressed.
\formula{$\text{EVA}_t = \text{NOPAT}_t - \text{WACC}\times \text{IC}_{t-1}$}
\formula{$EV = \text{IC}_0 + \sum \dfrac{\text{EVA}_t}{(1+\text{WACC})^t}$}
The identity holds only if (i) FCFF is defined consistently as
$\text{NOPAT}-\Delta\text{IC}$, so that capital charged equals capital invested, and
(ii) the terminal year is a genuine steady state, so the perpetuity of EVA matches the
perpetuity of FCFF. Drop either and the two answers diverge.
\end{soln}
\ans{b}{FCFF $=$ NOPAT $-\,\Delta$IC, and a true steady-state terminal year}

\q{17}{When methods disagree}{Your FCFF valuation gives Rs 520 per share and your FCFE
valuation gives Rs 480. What should you do?}
\opt{Average them to Rs 500}{Report the more conservative figure}{Find and fix the
inconsistency}{Report both as a range}
\begin{soln}
The methods are algebraically identical when the five consistencies hold. A gap is not
a modelling choice or a legitimate range --- it is proof of an error.\par
The three usual culprits, in order of frequency: debt not growing at $g$ where a
constant WACC was assumed; the tax shield double-counted (in the rate \emph{and} in the
flow); cash or investments treated differently in the flows and in the bridge.
\end{soln}
\ans{c}{Find and fix the inconsistency}
\memo{Never average. Averaging two answers that must be equal hides the error and
throws away the diagnostic.}

% ================================================================= SECTION C
\sect{C}{Cost of capital, relevering and circularity}{Nine questions on the inputs that
decide the discount rate --- and on the two places where students most often break risk
consistency without noticing.}

\q{18}{Risk-free rate}{The 10-year G-Sec yield is 6.85\% and the India sovereign default
spread is 1.87\%. The equity risk premium being used already includes a country risk
premium. What risk-free rate should be used?}
\opt{6.85\%}{8.72\%}{4.98\%}{7.08\%}
\begin{soln}
The G-Sec is not default-free: India is rated Baa3, so the yield contains compensation
for sovereign default risk. If the ERP already carries a country risk premium, using
the raw G-Sec counts India risk \emph{twice}.
\formula{$r_f = 6.85\% - 1.87\% = \mathbf{4.98\%}$}
The rule is symmetric: strip the spread from the risk-free rate \emph{or} use an ERP
without a country premium. Do one, never both and never neither.
\end{soln}
\ans{c}{4.98\%}
\memo{Double-counting country risk is a silent error --- the number still looks
plausible. State which convention you used.}

\q{19}{Relevering --- Hamada}{Unlevered beta is 0.90, market D/E is 6.4681\% and the
marginal tax rate is 25.168\%. What is the levered beta under the fixed-debt (Hamada)
formula?}
\opt{0.9436}{0.9582}{0.9729}{0.9000}
\begin{soln}
\formula{$\beta_L = \beta_U\left[1+(1-t)\tfrac{D}{E}\right]$}
$= 0.90\,[1 + (1-0.25168)(0.064681)]$\par
$= 0.90\,[1 + (0.74832)(0.064681)] = 0.90 \times 1.048403 = \mathbf{0.9436}$.
\end{soln}
\ans{a}{0.9436}
\memo{The tax term is $(1-t)$, not $(1+t)$. A sign slip here is a real error that
appeared in one of the project workbooks.}

\q{20}{Relevering --- Harris--Pringle}{The same company states a policy of constant
debt-to-value with continuous rebalancing, and discounts its tax shield at $\rho$. What
is the correct levered beta?}
\opt{0.9436}{0.9582}{0.9000}{0.9729}
\begin{soln}
Under Case 2 --- constant $D/V$, continuously rebalanced --- the debt level moves with
firm value, so the tax shield carries \emph{asset} risk rather than debt risk. It is
discounted at $\rho$, and the relevering formula carries \textbf{no tax term}.
\formula{$\beta_L = \beta_U\left(1+\tfrac{D}{E}\right) = 0.90\,(1.064681)
= \mathbf{0.9582}$}
Using the Hamada formula while discounting the shield at $\rho$ mixes Case 1 and Case 2
and breaks risk consistency.
\end{soln}
\ans{b}{0.9582}
\memo{Match the unlever, the relever and the shield discount rate to \emph{one} debt
policy. This is the single most common risk-consistency failure in real models.}

\q{21}{The $\rho$ identity}{WACC is 11.4286\%, $K_d$ is 8.34\%, the tax rate is 25.168\%
and $D/V$ is 6.0751\%. What is the unlevered cost of capital $\rho$?}
\opt{11.4286\%}{11.5561\%}{11.3011\%}{11.7641\%}
\begin{soln}
\formula{$\rho = \text{WACC} + K_d\,t\,\tfrac{D}{V}$}
$= 0.114286 + (0.0834)(0.25168)(0.060751)$\par
$= 0.114286 + 0.001275 = \mathbf{11.5561\%}$.\par
Equivalently $\rho = \tfrac{E}{V}K_e + \tfrac{D}{V}K_d$ --- the unweighted mix with no
tax relief, because under CCF and APV the shield has been taken out of the rate.
\end{soln}
\ans{b}{11.5561\%}
\memo{This identity is the fastest audit in the whole model. If $\rho$ and WACC do not
satisfy it to the last decimal, something is wrong.}

\q{22}{Debt policy and the shield}{Under which debt policy is the interest tax shield
discounted at the cost of debt $K_d$?}
\opt{Constant $D/V$, continuously rebalanced}{Debt fixed in rupee terms}{Constant
$D/V$, rebalanced annually}{All three}
\begin{soln}
The discount rate for the shield must match the risk of the shield, which follows the
risk of the debt.
\formula{Case 1 --- debt fixed in Rs: shield is as safe as the debt $\to$ discount at $K_d$}
\formula{Case 2 --- Harris--Pringle: debt moves with value $\to$ discount at $\rho$}
\formula{Case 3 --- Miles--Ezzell: first year at $K_d$, later years at $\rho$}
Only fixed rupee debt gives a shield whose amount is known in advance and therefore
carries debt risk throughout.
\end{soln}
\ans{b}{Debt fixed in rupee terms}

\q{23}{Circularity}{WACC needs market-value weights, but the equity value is what you
are solving for. What are the two sanctioned solutions?}
\opt{Use book-value weights, or ignore the problem}{Use a target D/E, or iterate to
convergence}{Use the industry average, or use book values}{Use the risk-free rate
instead of WACC}
\begin{soln}
Both solutions are legitimate and both must be disclosed.\par
\textbf{(i) Target D/E.} Use management's stated or the industry's structure.
Appropriate when the firm will migrate to it. Only the \emph{ratio} is needed, never
rupee values.\par
\textbf{(ii) Iterative convergence.} Seed with market capitalisation, compute
WACC $\to$ value $\to$ implied equity, feed that back and repeat until
$(\text{implied} - \text{assumed}) = 0$. In Excel: enable iterative calculation or use
Goal Seek.\par
For a listed company, using market cap is acceptable \emph{if stated} --- but note the
tension: you are using the very price you claim may be wrong.
\end{soln}
\ans{b}{Use a target D/E, or iterate to convergence}
\memo{Never silently use book-value weights. That is an automatic mark lost.}

\q{24}{Tax rate}{A company has migrated to section 115BAA. What is the effective
marginal corporate tax rate?}
\opt{22.000\%}{25.168\%}{34.944\%}{17.160\%}
\begin{soln}
Section 115BAA is a 22\% base rate, plus a 10\% surcharge, plus a 4\% health and
education cess. The loadings are multiplicative, not additive.
\formula{$22\% \times 1.10 \times 1.04 = \mathbf{25.168\%}$}
MAT does not apply under 115BAA. The old regime works out near 34.944\% and section
115BAB (new manufacturing) near 17.16\%. Always verify the regime from the tax
reconciliation note --- never assume.
\end{soln}
\ans{b}{25.168\%}

\q{25}{Cost of debt}{Which is the preferred source for the pre-tax cost of debt?}
\opt{Interest expense divided by average debt}{The yield on the company's traded bonds,
or the rate disclosed in the borrowings note}{The risk-free rate plus 2\%}{The coupon
on the oldest outstanding loan}
\begin{soln}
$K_d$ must be a \emph{forward-looking market} rate --- what the company would pay to
borrow today. Interest expense over average debt is a historical average of legacy
borrowings at legacy rates, distorted by drawdowns and repayments during the year.\par
Use it only as a sanity check. And remember the after-tax conversion for the WACC:
\formula{$K_d^{\text{after-tax}} = K_d(1-t)$}
\end{soln}
\ans{b}{The yield on the company's traded bonds, or the rate disclosed in the
borrowings note}

\q{26}{Beta}{A regression beta from 60 monthly observations is 1.15 with an $R^2$ of
0.27. The bottom-up beta from 469 sector peers is 0.91. Which should you use, and why?}
\opt{The regression beta, because it uses the company's own price history}{The
bottom-up beta, because the regression explains little of the variance}{The average of
the two}{Whichever gives the higher valuation}
\begin{soln}
An $R^2$ of 0.27 means the market explains only 27\% of the stock's variation --- the
regression estimate carries a wide standard error and is unstable to the choice of
window and frequency.\par
A bottom-up beta averages the unlevered betas of many firms in the same business, so
estimation errors diversify away, and it is then relevered at the subject's own D/E.
Where the two diverge materially, prefer bottom-up and \emph{explain} the choice.
\end{soln}
\ans{b}{The bottom-up beta, because the regression explains little of the variance}
\memo{Note the direction: here bottom-up is lower, so it \emph{raises} the valuation.
Say so --- disclosing that your choice helped the answer is what earns the mark.}

% ================================================================= SECTION D
\sect{D}{Terminal value, multiples and relative valuation}{Terminal value usually
carries 70--80\% of enterprise value, so these are the highest-leverage assumptions in
any model --- and the ones examiners probe hardest.}

\q{27}{Nominal growth}{Long-run real GDP growth is 6.5\% and long-run inflation 4.0\%.
What is nominal GDP growth, built the way the method requires?}
\opt{10.50\%}{10.76\%}{2.50\%}{26.00\%}
\begin{soln}
Growth rates compound; they do not add.
\formula{$g_{\text{nom}} = (1+g_{\text{real}})(1+i)-1 = (1.065)(1.040)-1 = \mathbf{10.76\%}$}
The additive shortcut gives 10.50\%, which is 26 basis points too low. Over a
perpetuity that error is not trivial.
\end{soln}
\ans{b}{10.76\%}
\memo{Terminal $g$ must also stay below WACC, and a mature single-country firm should
not out-grow nominal GDP forever.}

\q{28}{Steady state}{Terminal growth is 6.99\% and terminal ROIC is 11.94\%. What
reinvestment rate does steady state require?}
\opt{41.4\%}{58.6\%}{69.9\%}{11.9\%}
\begin{soln}
In steady state, growth is funded entirely by reinvesting a fraction of NOPAT at the
return on new capital.
\formula{$g = \text{RONIC}\times\text{reinvestment rate} \;\Rightarrow\;
\text{reinvestment rate} = \dfrac{g}{\text{RONIC}}$}
$= 0.0699 / 0.1194 = \mathbf{58.6\%}$.\par
So terminal FCFF $=$ NOPAT $\times (1-0.586) = 41.4\%$ of NOPAT. Forgetting to charge
this reinvestment is the fastest way to overstate a terminal value.
\end{soln}
\ans{b}{58.6\%}

\q{29}{Justified multiple}{Terminal NOPAT/EBITDA is 0.4902, terminal ROIC 11.94\%,
terminal growth 6.99\% and $\rho$ 11.39\%. What EV/EBITDA do these fundamentals justify?}
\opt{19.4x}{4.65x}{7.52x}{1.13x}
\begin{soln}
\formula{$\dfrac{EV}{EBITDA} = \dfrac{(NOPAT/EBITDA)\left(1-g/ROIC\right)}{WACC-g}$}
Reinvestment rate $= 0.0699/0.1194 = 0.5858$, so $1-0.5858 = 0.4142$.\par
Numerator $= 0.4902 \times 0.4142 = 0.2031$.\par
Denominator $= 0.1139 - 0.0699 = 0.0440$.\par
$EV/EBITDA = 0.2031/0.0440 = \mathbf{4.65x}$.
\end{soln}
\ans{b}{4.65x}
\memo{The test is not whether your multiple equals the peer multiple. It is whether
your multiple is arithmetically consistent with your own $g$, ROIC and WACC.}

\q{30}{EV to invested capital}{A company earns ROIC of 3.3\% today, terminal growth is
6.99\% and WACC 11.36\%. What multiple of invested capital is justified?}
\opt{1.66x}{0.46x}{Negative --- no positive multiple is justified}{2.02x}
\begin{soln}
\formula{$\dfrac{EV}{IC} = \dfrac{ROIC-g}{WACC-g}
= \dfrac{0.033-0.0699}{0.1136-0.0699} = \dfrac{-0.0369}{0.0437} = \mathbf{-0.85x}$}
The numerator is negative because the company earns \emph{less} than the rate at which
it is being asked to grow. Every rupee reinvested destroys value, so no positive
multiple of invested capital is warranted at all.\par
This is the sharpest single statement of the Ambuja investment case.
\end{soln}
\ans{c}{Negative --- no positive multiple is justified}
\memo{When ROIC $<g$, growth is value-destroying. When ROIC $=$ WACC, growth is
value-neutral. Only ROIC $>$ WACC makes growth worth having.}

\q{31}{Reverse cross-check}{A peer median exit multiple implies terminal growth of
10.33\%, while nominal GDP growth is 10.76\% and the company's realisation has
compounded at 2.0\% over ten years. What is the correct conclusion?}
\opt{Adopt the peer multiple --- the market knows best}{The peer multiple is the
implausible end of the range, and you should say so}{Raise your terminal growth to
10.33\%}{Relative valuation cannot be reconciled with DCF, so ignore it}
\begin{soln}
Solving a peer multiple back for the growth it implies is the reverse cross-check the
method requires. Here it demands perpetual growth at roughly 96\% of the whole
economy's rate, from a company whose price has compounded at 2.0\%.\par
A company growing faster than the economy forever eventually \emph{becomes} the
economy. Adopting that growth to justify the multiple would contradict your own model.
Take a position, evidence it, and record the gap honestly.
\end{soln}
\ans{b}{The peer multiple is the implausible end of the range, and you should say so}

\q{32}{Terminal dominance}{Two models value similar companies. Model A uses a five-year
explicit horizon and terminal value is 79\% of EV; Model B uses ten years and terminal
value is 72\%. What does this mainly reflect?}
\opt{Model B is more optimistic}{Horizon length, not optimism}{Model A has a higher
WACC}{Model B has higher terminal growth}
\begin{soln}
Lengthening the explicit forecast moves cash flows out of the terminal perpetuity and
into the modelled period. The \emph{same} company modelled over ten years will always
show a lower terminal share than over five.\par
The method's guidance: extend the explicit period until steady state is credible ---
particularly when ROIC $\gg$ WACC or growth is far above steady state. Report terminal
value as a percentage of EV, and treat 80--90\% as the level demanding scrutiny.
\end{soln}
\ans{b}{Horizon length, not optimism}

\q{33}{Choice of multiple}{For which type of business is price-to-book the appropriate
primary multiple?}
\opt{FMCG and IT services}{Banks, NBFCs, oil and gas, real estate}{Any company with
negative earnings}{Companies with high operating leverage}
\begin{soln}
The selection rule follows where value actually sits on the financial statements.
\formula{Income-statement businesses $\to$ P/E or EV/EBITDA}
\formula{Balance-sheet businesses $\to$ P/B, paired with ROE}
For a bank, the balance sheet \emph{is} the business, so book value is the meaningful
scale variable. P/E is meaningless with negative earnings; P/B needs positive net
worth.
\end{soln}
\ans{b}{Banks, NBFCs, oil and gas, real estate}

\q{34}{Valuing a bank}{Why is an enterprise FCF model inappropriate for a bank, and
what should be used instead?}
\opt{Banks have no cash flows; use EV/EBITDA}{Debt and working capital are the raw
material, so value equity directly with DDM or FCFE at $K_e$}{Banks are too volatile;
use book value only}{Use FCFF but with a higher WACC}
\begin{soln}
For an industrial firm, debt is financing. For a bank, deposits and borrowings are
\emph{inputs to production} --- the thing it buys and sells. Separating operating from
financing flows is therefore meaningless, and enterprise value has no clear content.
\formula{Value equity directly: DDM or FCFE, discounted at $K_e$}
Reinvestment is driven by capital-adequacy requirements rather than by capex. The
primary multiple is P/B paired with ROE, and residual income works well too.
\end{soln}
\ans{b}{Value equity directly --- DDM or FCFE at $K_e$}
\memo{Stating this method switch prominently is itself graded --- it demonstrates you
understand risk consistency rather than just following a template.}

\q{35}{Order of work}{Where does relative valuation belong in the workflow, and why?}
\opt{Before the DCF, to anchor the target price}{After the DCF, because a multiple is
an outcome of valuation rather than its cause}{In place of the DCF when time is
short}{Only for loss-making companies}
\begin{soln}
A multiple is value divided by an accounting indicator --- it is the \emph{result} of a
valuation, not an input to one. Anchoring a DCF to ``the market multiple says so''
reverses the logic and imports the market's assumptions into what is meant to be an
independent estimate.\par
The correct sequence: build the DCF, derive the multiple your own model implies, then
place it beside the peer set and \emph{explain} any gap through fundamentals ---
growth for P/E, ROE for P/B, returns for EV/IC.
\end{soln}
\ans{b}{After the DCF, because a multiple is an outcome of valuation rather than its
cause}

% ================================================================= SECTION E
\sect{E}{Applied: the cement valuations}{Five questions using the actual figures from
the Ambuja and UltraTech models, where the theory above has to survive contact with
real accounts.}

\q{36}{Basis error}{Ambuja's group market capitalisation is Rs 99,095 cr. Standalone
EBITDA is Rs 2,930 cr; consolidated EBITDA is Rs 6,559 cr; non-controlling interest is
Rs 12,499 cr and the group is roughly net-cash. Which EV/EBITDA is correct, and why?}
\opt{34.0x, on standalone EBITDA}{17.0x, on consolidated EBITDA with NCI added to
EV}{15.1x, on consolidated EBITDA without NCI}{Either is acceptable if disclosed}
\begin{soln}
Market capitalisation prices the \emph{whole group}, including the 50.05\% stake in the
separately listed ACC. Standalone EBITDA excludes ACC entirely. Dividing one by the
other is a category error.
\formula{Wrong: $99{,}676 / 2{,}930 = 34.0x$}
\formula{Right: $(99{,}095 + 12{,}499 - 93)/6{,}559 = 111{,}501/6{,}559 = \mathbf{17.0x}$}
Non-controlling interest is added to enterprise value because consolidated EBITDA
includes 100\% of the partly owned subsidiaries that produce it. Numerator and
denominator must cover the same set of assets.
\end{soln}
\ans{b}{17.0x, on consolidated EBITDA with NCI added to EV}
\memo{Whenever a multiple looks absurd, check the basis before you check the market.}

\q{37}{The bridge}{Ambuja's consolidated cash flows include 100\% of its partly owned
subsidiaries. Non-controlling interest is Rs 12,499 cr and there are 247.18 cr shares.
What happens if the NCI is not deducted in the bridge?}
\opt{Nothing --- NCI is a book entry}{Value per share is overstated by about
Rs 51}{Value per share is understated by about Rs 51}{Only the WACC changes}
\begin{soln}
If the cash flows being discounted belong 100\% to the group, but only part of the
group belongs to Ambuja's shareholders, the minority's share must be removed before
dividing by Ambuja's share count.
\formula{$12{,}499 / 247.18 = \mathbf{Rs\ 50.57}$ per share}
Omitting it hands the parent's shareholders value that belongs to the minorities ---
liability consistency again, in its most expensive form.
\end{soln}
\ans{b}{Value per share is overstated by about Rs 51}

\q{38}{After-tax returns}{A workbook computes historical ROIC as EBIT of Rs 11,761 cr
over invested capital of Rs 81,733 cr, but the cell meant to apply tax references an
empty cell. Statutory tax is 25.168\%. What is the reported ROIC, and what is the
correct one?}
\opt{Reported 14.39\%, correct 10.77\%}{Reported 10.77\%, correct 14.39\%}{Both 14.39\%}{Reported 14.39\%, correct 12.50\%}
\begin{soln}
With the tax reference broken, NOPAT collapses to EBIT and the series becomes a
\emph{pre-tax} return.
\formula{Reported: $11{,}761/81{,}733 = 14.39\%$ \quad (pre-tax)}
\formula{Correct: $11{,}761 \times (1-0.25168)/81{,}733 = \mathbf{10.77\%}$}
The distinction matters enormously here: 14.39\% sits comfortably above an 11.4\% cost
of capital, while 10.77\% sits \emph{below} it. The same company looks value-creating
or value-destroying depending on one cell reference.
\end{soln}
\ans{a}{Reported 14.39\%, correct 10.77\%}
\memo{ROIC is always after tax. A pre-tax return compared with a WACC is comparing
unlike with unlike.}

\q{39}{Normalising a tax credit}{Ambuja reports a net tax \emph{credit} of Rs 2,338 cr,
so reported profit exceeds pre-tax profit and the effective rate is minus 71\%. How
should the valuation treat it?}
\opt{Use the reported effective rate of $-71\%$}{Use the statutory rate throughout and
decompose the credit by its economics}{Ignore tax entirely for that year}{Average the
reported rate over five years}
\begin{soln}
A negative effective rate is unusable as a forward-looking assumption. Tax EBIT at the
statutory 25.168\% throughout, then treat each component of the credit on its
economics, not its accounting:\par
$\bullet$ Provision reversals (Rs 1,251 cr) --- non-cash, prior period. \textbf{Excluded}:
the credit was never in the cash flows.\par
$\bullet$ Tax paid under protest (Rs 614 cr) --- real cash already handed over,
recoverable. \textbf{Added} as a non-operating asset.\par
$\bullet$ Deferred tax liability (Rs 3,466 cr) --- a real claim. \textbf{Deducted} once
in the bridge.\par
Net effect: about \textbf{minus Rs 11.54 per share}. A Rs 2,338 cr headline gain is
worth negative value.
\end{soln}
\ans{b}{Use the statutory rate throughout and decompose the credit by its economics}

\q{40}{EVA and enterprise value}{A company's EVA is negative in nine of ten forecast
years. What must be true of its enterprise value?}
\opt{It must be negative}{It must sit below invested capital}{It must exceed invested
capital}{Nothing follows}
\begin{soln}
\formula{$EV = \text{IC}_0 + \text{PV(future EVA)}$}
If EVA is persistently negative, the present value of future EVA is negative, so EV
falls below opening invested capital. That is not an error --- it is exactly what the
theory predicts when a company earns less on capital than capital costs.\par
EV can still be comfortably positive; it is simply worth less than the capital sunk
into it. Equivalently, EV/IC $<1$, which the identity in Q30 also produces.
\end{soln}
\ans{b}{It must sit below invested capital}
\memo{Negative EVA with a positive EV is perfectly coherent. Do not ``fix'' it.}

% ================================================================= SECTION F
\sect{F}{Guest lecture --- M\&A, synergy and exchange ratios}{Five questions on the
transaction module. Value each entity standalone first, then value only the
\emph{incremental} synergy cash flows, at a rate that matches their risk.}

\q{41}{Definition of synergy}{Acquirer A is worth Rs 800 cr standalone and target B
Rs 300 cr standalone. The combined entity is valued at Rs 1,250 cr. What is the synergy
value?}
\opt{Rs 1,250 cr}{Rs 450 cr}{Rs 150 cr}{Rs 950 cr}
\begin{soln}
Synergy is the value the combination creates over and above what the two businesses are
already worth apart.
\formula{$\text{Synergy} = V_{AB} - (V_A + V_B) = 1{,}250 - (800+300) = \mathbf{Rs\ 150\ cr}$}
Note what this requires: both standalone valuations must be built \emph{first} and on a
consistent basis. Synergy is a residual, so any error in either standalone figure lands
entirely in the synergy number.
\end{soln}
\ans{c}{Rs 150 cr}

\q{42}{Net gain to the acquirer}{On the same facts, A pays Rs 380 cr for B. What is the
net gain to A's shareholders?}
\opt{Rs 150 cr}{Rs 80 cr}{Rs 70 cr}{Rs 380 cr}
\begin{soln}
The premium is what the acquirer hands to the target's shareholders; the net gain is
what is left of the synergy after paying it.
\formula{Premium $= 380 - 300 = \text{Rs }80$ cr}
\formula{Net gain to A $=$ Synergy $-$ Premium $= 150 - 80 = \mathbf{Rs\ 70\ cr}$}
The transaction creates Rs 150 cr of value, of which A's shareholders keep Rs 70 cr and
B's shareholders capture Rs 80 cr. If the premium had exceeded Rs 150 cr, A would be
transferring value to B despite a genuinely synergistic deal.
\end{soln}
\ans{c}{Rs 70 cr}
\memo{A good deal at a bad price is a bad deal. Always report the net gain after the
premium, not the synergy alone.}

\q{43}{Exchange ratio}{A share-swap merger values the target at Rs 480 per share and the
acquirer at Rs 1,600 per share. What is the exchange ratio?}
\opt{0.30 acquirer shares per target}{3.33 acquirer shares per target}{0.48 acquirer shares per target}{1.60 acquirer shares per target}
\begin{soln}
The exchange ratio equates the value given to the value received, per share.
\formula{$\text{ER} = \dfrac{\text{value per target share}}{\text{value per acquirer share}}
= \dfrac{480}{1{,}600} = \mathbf{0.30}$}
So each target shareholder receives 0.30 acquirer shares. If the target has 100 cr
shares, the acquirer issues 30 cr new shares.\par
The ratio should be set on \emph{intrinsic} values, including the target's share of
synergy that the acquirer is willing to pay away --- not on traded prices alone, which
already contain deal speculation.
\end{soln}
\ans{a}{0.30 acquirer shares per target share}

\q{44}{EPS accretion}{A merger is EPS-accretive in year one. What does this tell you
about value creation?}
\opt{The deal creates value}{Nothing reliable --- accretion depends on the relative P/E
ratios}{The deal destroys value}{The synergies must be real}
\begin{soln}
A company with a \emph{higher} P/E acquiring one with a \emph{lower} P/E will report
EPS accretion arithmetically, with no synergy whatever, simply because it buys more
earnings per share issued. The reverse produces dilution even when the deal is
excellent.
\formula{Accretion is an artefact of relative P/E, not evidence of economic value}
Value creation is tested by discounting incremental synergy cash flows and comparing
them with the premium paid --- Q41 and Q42, not the EPS line.
\end{soln}
\ans{b}{Nothing reliable --- accretion depends on the relative P/E ratios}
\memo{Never substitute EPS accretion for economic value creation. This is the single
most tested point in the transaction module.}

\q{45}{Valuing synergies}{Which set of cash flows may be included when valuing synergy,
and at what rate?}
\opt{All the target's cash flows, at the acquirer's WACC}{Only incremental, timing-specific flows, at a risk-matched rate}{The combined entity's total cash flows, at the target's WACC}{Cost savings only, at the risk-free rate}
\begin{soln}
Three disciplines apply.\par
\textbf{Incremental.} Only cash flows that would not exist without the combination.
The target's standalone flows are already in $V_B$; counting them again double-counts.\par
\textbf{Timing-specific.} Synergies arrive on a schedule and cost money to achieve.
Implementation costs and their timing must be explicit.\par
\textbf{Risk-matched.} Cost synergies are usually more certain than revenue synergies,
so they do not deserve the same discount rate. Revenue synergies should carry a higher
rate, or a probability weighting.\par
Treat revenue, cost, tax and financing effects separately, and sensitise the
assumptions that drive the premium decision.
\end{soln}
\ans{b}{Only incremental, timing-specific cash flows, discounted at a rate matching
their own risk}
\memo{Revenue synergies are the ones that most often fail to arrive. Examiners expect
you to say so.}

\vspace{8mm}
\begin{center}
{\color{inkred}\bfseries\large --- end of question bank ---}\par\vspace{2mm}
{\color{pencil}\itshape 45 questions \quad\textbar\quad 40 core $+$ 5 guest lecture}
\end{center}

\end{document}
